\begin{table*}[t]
\centering
\caption{Total-cost-of-ownership scenario: brackish-water pump impeller, 25-year horizon, zero discount rate. Costs \prop{} except anchors in Table~\ref{tab:costmodel}. Bronze baseline: \$2.5k purchase, 6-yr interval (5 purchases; 4 downtime events). WAAM-HEA: \$3k first cost, interval $T$. $D$ = downtime-plus-labor per replacement event.}
\label{tab:tcobreakeven}
\footnotesize
\begin{tabularx}{\textwidth}{@{}lXXXX@{}}
\toprule
Downtime $D$ per event & Bronze TCO (25~yr) & HEA TCO at $T=6$~yr & HEA TCO at $T=12$~yr & Breakeven rule for HEA \\
\midrule
\$8k (critical service) & $12.5+4\times8 = \$44.5$k & $3+3\times11 = \$36$k & $3+1\times11 = \$14$k & wins if $T \geq \sim$5~yr (below bronze interval) \\
\$2k (typical berth) & $12.5+4\times2 = \$20.5$k & $3+3\times5 = \$18$k & $3+1\times5 = \$8$k & wins if $T \geq$ 6~yr (match bronze) \\
\$0.5k (non-critical) & $12.5+4\times0.5 = \$14.5$k & $3+3\times3.5 = \$13.5$k & $3+1\times3.5 = \$6.5$k & wins if $T \geq \sim$6~yr; margin thin \\
\bottomrule
\end{tabularx}

\vspace{4pt}
\footnotesize Scaling rule \prop{}: required HEA interval $\approx$ incumbent interval when downtime $\gtrsim$ \$1--2k/event; $\approx$1.2--1.3$\times$ incumbent when downtime is negligible. LCA: each avoided replacement amortizes the Co-bearing feedstock burden and atomization energy; fouling-release benefit of electropolished HEA surfaces is plausible but unproven \cite{nass2025marine}.

\end{table*}
